% SHARD-ID: CA-25-GAUSSIAN-BOSON-DIAGONALIZATION
% SHARD-TITLE: Gaussian Boson Diagonalization and the One-Particle Symbol
% SHARD-SUMMARY: Diagonalizes the scalar translation-invariant Gaussian boson tier into oscillator modes in spatial dimensions d=1,2,3.
% SHARD-SUMMARY: Separates the sourced free-scalar Fock construction from the local coefficient-symbol derivation used by the compiler.
% SHARD-SUMMARY: Records the exact boundary where general bosonic BdG and pairing systems remain deferred.
% SHARD-KEYWORDS: Gaussian boson, diagonalization, one-particle Hilbert space, Fock, dispersion, BdG deferral

\section{Gaussian Boson Diagonalization and the One-Particle Symbol}
\label{sec:gaussian-boson-diagonalization}

\subsection{Status, Sources, and Dependencies}

\textbf{Status: sourced template plus flat local coefficient derivation.}  The
standard lattice free-scalar Hamiltonian, Fock representation, continuum
one-particle space, and second-quantized Hamiltonian are sourced.  The extension
from nearest-neighbour coefficients to an arbitrary real translation-invariant
scalar symbol is a local algebraic calculation on the flat \(L^2(dk)\) symbol
space fixed in CONVENTIONS.md (j).

Dependencies: CA-23--CA-24 and CONVENTIONS.md (a), (f), (j).

% Source: literature/md/2010.11121/2010.11121.md:251 -- Weyl algebra and one-particle bosonic Fock representation.
% Source: literature/md/2010.11121/2010.11121.md:274 -- hypercubic lattice and discrete Fourier setup.
% Source: literature/md/2010.11121/2010.11121.md:598 -- sourced free lattice Hamiltonian.
% Source: literature/md/2010.11121/2010.11121.md:603 -- sourced lattice dispersion.
% Source: literature/md/2010.11121/2010.11121.md:648--672 -- continuum OAR one-particle scalar product and its flat Fock realization.
% Source: literature/md/2010.11121/2010.11121.md:670 -- continuum Fock/Weyl representation by creation and annihilation operators.
% Source: literature/md/2010.11121/2010.11121.md:1037 -- continuum Hamiltonian as second quantization of the one-particle generator.
% Source: literature/md/2010.11121/2010.11121.md:1540 -- one-particle Hamiltonians as multiplication by dispersion in a cutoff diagonalization.
% Source: literature/md/2010.11121/2010.11121.md:1626--1628 -- massless fields require a zero-mode policy.
% Source: references/cft/Schottenloher2008/Schottenloher2008.md:4186 -- free bosonic QFT example assumes m > 0.

\subsection{Scalar Coefficient Tier}

Fix \(d\in\{1,2,3\}\).  The first Gaussian tier is
\[
  H
  =
  \frac{1}{2}\sum_x p_x^2
  +
  \frac{1}{2}\sum_{x,r}q_xV_rq_{x+r},
  \qquad
  V_r=V_{-r}\in\mathbb R .
\]
On an infinite translation-invariant lattice, or on a periodic lattice before
taking volume limits, the Fourier transform sends the potential kernel to
\[
  \omega_\epsilon(k)^2
  =
  \sum_r V_r e^{i\epsilon k\cdot r}.
\]
The reality condition \(V_r=V_{-r}\) makes this symbol real:
\[
  \omega_\epsilon(k)^2
  =
  V_0+\sum_{r>0}2V_r\cos(\epsilon k\cdot r).
\]
The stability condition for this tier is
\[
  \omega_\epsilon(k)^2\ge 0
\]
on the Brillouin zone.  The \texttt{mass} parameter is a nonnegative label; in
the nearest-neighbour Klein-Gordon family, \(m>0\) gives the positive-dispersion
case used by the current generator statements.  For a general coefficient
symbol the required condition remains \(\omega_\epsilon(k)>0\) on the chosen
patch.  When \(m=0\), zero modes are coefficient data only until a zero-mode
convention excludes or otherwise controls them.

\subsection{Mode Decoupling}

The Hamilton equations are a local derivation from the canonical commutators:
\[
  \dot q_x=p_x,
  \qquad
  \dot p_x=-\sum_rV_rq_{x+r}.
\]
Fourier transforming gives
\[
  \dot{\widehat q}(k)=\widehat p(k),
  \qquad
  \dot{\widehat p}(k)=-\omega_\epsilon(k)^2\widehat q(k),
\]
and hence
\[
  \ddot{\widehat q}(k)+\omega_\epsilon(k)^2\widehat q(k)=0.
\]
Thus each momentum mode is an oscillator with frequency \(\omega_\epsilon(k)\).
This is the algebraic reason the compiler can reduce the scalar Gaussian tier
to multiplication by a dispersion symbol.

\subsection{Complex Structure and Fock Space}

For \(\omega>0\) on the chosen momentum patch, the oscillator complex structure is
\[
  (q,p)\longmapsto (-\omega p,\omega^{-1}q).
\]
The sourced free-scalar construction uses this pattern with
\(\gamma_m(k)=\sqrt{|k|^2+m^2}\) in the continuum and with
\(\gamma_m^{(N)}(k)\) on the lattice.  It also realizes the Weyl algebra in a
bosonic Fock representation by creation and annihilation operators.
The source scalar product is the OAR one on real Cauchy data, with
\(\gamma_m^{-1/2}\widehat q+i\gamma_m^{1/2}\widehat p\); the text then realizes
the representation on a mass-independent flat \(L^2\) Fock space.  This shard
uses only the resulting flat symbol bookkeeping and does not prove that the
boost formula of CA-24/CA-26 intertwines the OAR or mass-shell representation.

The compiler-level output for the positive-dispersion scalar tier is therefore:
\[
\begin{array}{ll}
\mathfrak h_\epsilon
  & \text{flat one-particle momentum space on the chosen patch or grid},\\
h_\epsilon
  & \text{multiplication by }\omega_\epsilon(k),\\
\mathcal F_+(\mathfrak h_\epsilon)
  & \text{symmetric Fock completion},\\
H_\epsilon
  & d\Gamma(h_\epsilon) + E_{\epsilon,0}\text{ if a vacuum-energy policy is named}.
\end{array}
\]
The additive constant \(E_{\epsilon,0}\) is not part of the commutator algebra
for the generators considered here, but it is part of a Hamiltonian-density
normal-ordering policy and must be recorded before using ground-state energies
as numerical claims.

\subsection{What Is Not Yet Included}

This shard does not define the full bosonic BdG tier.  A general quadratic
Hamiltonian in creation and annihilation operators may include pairing terms,
and then the correct diagonalization is a symplectic/Bogoliubov problem with
positivity, implementability, and domain conditions.  Those conventions are not
fixed here.

The current compiler rule is deliberately narrower:
\[
  \text{scalar }(q,p)\text{ tier}
  \quad\Rightarrow\quad
  \omega_\epsilon(k)^2\text{ symbol}
  \quad\Rightarrow\quad
  h_\epsilon=\omega_\epsilon(k).
\]
BdG inputs must fail closed or route to a later convention until that tier is
written.

\subsection{Compiler Consequence}

For the positive-dispersion scalar Gaussian block, later candidate generators can
be constructed first on the one-particle smooth momentum core and then
second-quantized.  The
coefficient data enter only through \(\omega_\epsilon(k)^2\) and its
derivatives, while the continuum proof still requires the CA-15 payload:
scaling maps, state convergence, domains, and observable covariance.
